\section{Class Relationships and UML Notation}
\label{sec:lld-relationships}

\problemstatement{Distinguish the ways two classes can relate --
association, aggregation, composition, dependency, and inheritance -- and
know the UML class-diagram notation, since interviewers expect you to
sketch a design and name these relationships correctly.}

\begin{patternbox}[Which Relationship Is It?]
\begin{itemize}
  \item \textbf{``is-a''} -- \textbf{inheritance} (a Car is-a Vehicle).
        UML: hollow triangle arrow to the parent.
  \item \textbf{``has-a,'' parts can live independently} --
        \textbf{aggregation} (a Team has Players; players outlive the
        team). UML: hollow diamond.
  \item \textbf{``has-a,'' parts die with the whole} --
        \textbf{composition} (a House has Rooms; destroy the house and
        the rooms go with it). UML: filled diamond.
  \item \textbf{``uses-a'' transiently} -- \textbf{dependency} (an
        \texttt{Order} uses a \texttt{PaymentGateway} passed to one
        method). UML: dashed arrow.
  \item \textbf{A general structural link} -- \textbf{association}
        (a plain solid line); aggregation and composition are its
        stronger, ownership-bearing forms.
\end{itemize}
\end{patternbox}

\subsection*{The key axis: ownership and lifetime}

The distinction interviewers probe is \textbf{aggregation vs.\
composition}, and it is entirely about \emph{lifetime}. In composition the
whole \emph{owns} its parts and controls their creation and destruction --
model it by storing parts \emph{by value} or via an owning
\texttt{unique\_ptr}. In aggregation the whole merely \emph{references}
parts that exist independently -- model it by storing a non-owning pointer
or reference handed in from outside.

\begin{ideabox}[Favour Composition Over Inheritance]
A foundational LLD maxim. Inheritance is rigid: it fixes behaviour at
compile time and leaks the base class's details into every subclass. Composing
an object out of smaller collaborators (holding a \texttt{Strategy}, an
\texttt{Engine}, a \texttt{Logger}) lets you change behaviour by swapping a
part at run time. When tempted to subclass, first ask whether ``has-a''
models the problem better than ``is-a.''
\end{ideabox}

\subsection*{C++ illustration of each relationship}

\begin{lstlisting}
#include <bits/stdc++.h>
using namespace std;

class Engine {                 // an independent part
public:
    void start() { cout << "Engine started\n"; }
};

class Logger {                 // dependency: used transiently
public:
    void log(const string& m) { cout << "[log] " << m << "\n"; }
};

// Composition: Car owns its Engine (created and destroyed with the Car).
class Car {
    Engine engine;                     // by value -> composition
public:
    void drive(Logger& logger) {       // dependency: Logger passed in
        engine.start();
        logger.log("car driving");
    }
};

// Aggregation: Team references Players it does not own.
class Player {
public:
    string name;
    Player(string n) : name(move(n)) {}
};
class Team {
    vector<Player*> players;           // non-owning -> aggregation
public:
    void add(Player* p) { players.push_back(p); }
    void roster() {
        for (auto* p : players) cout << p->name << " ";
        cout << "\n";
    }
};

int main() {
    Car car; Logger logger;
    car.drive(logger);

    Player a("Alice"), b("Bob");       // players outlive any team
    Team t; t.add(&a); t.add(&b);
    t.roster();
    return 0;
}
\end{lstlisting}

\begin{takeawaybox}
Name relationships precisely: inheritance is ``is-a''; composition is
owned ``has-a'' (filled diamond, shared lifetime); aggregation is
referenced ``has-a'' (hollow diamond, independent lifetime); dependency is
a transient ``uses-a.'' In code the tell is ownership -- by-value/owning
pointer for composition, non-owning pointer for aggregation. When in
doubt, prefer composition over inheritance.
\end{takeawaybox}
